% SEGMENT OLP-0469-S01
% Part: normal-modal-logic
% Chapter: tableaux
% Section: countermodels

\documentclass[../../../include/open-logic-section]{subfiles}

\begin{document}

\olfileid{nml}{tab}{cou}

\olsection{\usetoken{P}{tableau}களிலிருந்து எதிர்மாதிரிகள்}

\begin{explain}
  நிறைவுத்தன்மைத் தேற்றத்தின் நிறுவல், $\Entails !A$ எனில் $\Proves !A$
  என்று காட்டுவதோடு மட்டும் நிற்பதில்லை;~$\Entails/ !A$ எனில்
  % SOURCE CORRECTION TA-NML-054: restore the formula marker on A.
  $!A$ க்கான எதிர்மாதிரிகளைக் கட்டமைக்கும் ஒரு முறையையும் நமக்குத் தருகிறது.
  ~$\Log{K}$ இன் நேர்வில், இந்த முறை ஒரு \emph{தீர்மானச் செயல்முறையாக}
  அமைகிறது. $\Entails/ !A$ என்க. அப்போது
  \olref[cpl]{prop:complete-tableau} இன் நிறுவல், ஒரு முழுமையான
  !!a{tableau}வைக் கட்டமைக்கும் முறையைத் தருகிறது. உண்மையில், இந்த முறை
  எப்போதும் முடிவடைகிறது. $\Log{K}$ க்கான கூற்றுத் தருக்க விதிகள் குறைந்த
  சிக்கலுடைய முன்னொட்டுள்ள !!{formula}s ஐ மட்டுமே சேர்க்கின்றன; அதாவது,
  ஒவ்வொரு குறியிட்ட வாய்பாடு~$\sFmla{S}{!A}[\sigma]$ க்கும் ஒவ்வொரு
  கூற்றுத் தருக்க விதியையும் ஒரு கிளையில் ஒருமுறை மட்டுமே பயன்படுத்த வேண்டும்.
  புதிய முன்னொட்டுகள்
  \iftag{prvBox}{$\TRule{\False}{\Box}$}{}\iftag{notprvBox,notprvDiamond}{}{
    மற்றும் }\iftag{prvDiamond}{$\TRule{\True}{\Diamond}$}{}
  \iftag{notprvBox,notprvDiamond}{விதியால்}{விதிகளால்} மட்டுமே
  உருவாக்கப்படுகின்றன; அவற்றையும் ஒருமுறை மட்டுமே பயன்படுத்த வேண்டும் (மேலும்
  அவை ஒவ்வொன்றும் ஒரே ஒரு புதிய முன்னொட்டை உருவாக்கும்).
  \iftag{prvBox}{$\TRule{\True}{\Box}$}{}\iftag{notprvBox,notprvDiamond}{}{
    மற்றும் }\iftag{prvDiamond}{$\TRule{\False}{\Diamond}$}{}
  \iftag{notprvBox,notprvDiamond}{விதியை}{விதிகளை} பலமுறை பயன்படுத்த
  வேண்டியிருக்கலாம்; ஆனால் ஒரு முன்னொட்டிற்கு
  ஒருமுறை மட்டுமே, மேலும் முடிவுறு எண்ணிக்கையிலான புதிய முன்னொட்டுகளே
  உருவாக்கப்படுகின்றன. ஆகவே கட்டமைப்பு முடிவுறு எண்ணிக்கையிலான படிகளுக்குப்
  பிறகு ஒரு மூடிய கிளையையோ ஒரு முழுமையான கிளையையோ விளைவிக்கும்.

  திறந்த முழுமையான கிளையைக் கொண்ட ஓர் அட்டவணை மரம் கட்டமைக்கப்பட்டதும்,
  \olref[cpl]{thm:tableau-completeness} இன் நிறுவல், முன்னொட்டுள்ள
  !!{formula}s இன் மூலக் கணத்தை நிறைவுறுத்தும் வெளிப்படையான ஒரு மாதிரியைத்
  தருகிறது. ஆகவே $\Gamma \Entails !A$ எனில் ஒரு மூடிய !!{tableau} இருந்து
  $\Gamma \Proves !A$ என்பது மட்டுமல்ல; மூடிய !!{tableau}வை உரிய முறையில்
  தேடி, இறுதியில் ஒரு ``முழுமையான'' !!{tableau}வைப் பெற்றால்,
  $\Gamma \Entails/ !A$ என்று அறிந்துகொள்வதோடு உண்மையில் ஓர்
  எதிர்மாதிரியையும் கட்டமைக்க முடியும்.
\end{explain}

\iftag{prvBox}{
\begin{ex}
  $\Proves/ \Box(p \lor q) \lif (\Box p \lor \Box q)$ என்று அறிவோம்.
  ஓர் அட்டவணை மரத்தின் கட்டமைப்பு பின்வருமாறு தொடங்குகிறது:
  \begin{oltableau}
    [\pFmla{\False}{\Box(p \lor q) \lif (\Box p \lor \Box q)}{1},
      just = \TAss, checked
      [\pFmla{\True}{\Box(p \lor q)}{1},
        just = {\TRule{\False}{\lif}[1]},
        [\pFmla{\False}{\Box p \lor \Box q}{1},
          just = {\TRule{\False}{\lif}[1]}, checked
          [\pFmla{\False}{\Box p}{1},
            just = {\TRule{\False}{\lor}[3]}, checked
            [\pFmla{\False}{\Box q}{1},
              just = {\TRule{\False}{\lor}[3]}, checked
              [\pFmla{\False}{p}{1.1},
                just = {\TRule{\False}{\Box}[4]}, checked
                [\pFmla{\False}{q}{1.2},
                  just = {\TRule{\False}{\Box}[5]}, checked
                ]
              ]
            ]
          ]
        ]
      ]
    ]
  \end{oltableau}
  !!{tableau} இன்னும் முடிவடையவில்லை என்பது தெளிவு. அடுத்த படியில்,
  சரிபார்ப்புக்குறி இல்லாத ஒரே வரியை, அதாவது வரி~$2$ இலுள்ள முன்னொட்டுள்ள
  !!{formula}~$\sFmla{\True}{\Box(p \lor q)}[1]$ ஐக் கருதுகிறோம்.
  % SOURCE CORRECTION TA-NML-058: the search constructs a complete tableau, not a closed one.
  முழுமையான அட்டவணை மரத்தின் கட்டமைப்பின்படி, கிளையில் பயன்படுத்தப்பட்ட
  ஒவ்வொரு முன்னொட்டிற்கும், அதாவது $1.1$, $1.2$ ஆகிய இரண்டிற்கும்,
  $\TRule{\True}{\Box}$ விதியைப் பயன்படுத்த வேண்டும்:
  \begin{oltableau}
    [\pFmla{\False}{\Box(p \lor q) \lif (\Box p \lor \Box q)}{1},
      just = \TAss, checked
      [\pFmla{\True}{\Box(p \lor q)}{1},
        just = {\TRule{\False}{\lif}[1]},
        [\pFmla{\False}{\Box p \lor \Box q}{1},
          just = {\TRule{\False}{\lif}[1]}, checked
          [\pFmla{\False}{\Box p}{1},
            just = {\TRule{\False}{\lor}[3]}, checked
            [\pFmla{\False}{\Box q}{1},
              just = {\TRule{\False}{\lor}[3]}, checked
              [\pFmla{\False}{p}{1.1},
                just = {\TRule{\False}{\Box}[4]}, checked
                [\pFmla{\False}{q}{1.2},
                  just = {\TRule{\False}{\Box}[5]}, checked
                  [\pFmla{\True}{p \lor q}{1.1},
                    just = {\TRule{\True}{\Box}[2]}
                    [\pFmla{\True}{p \lor q}{1.2},
                      just = {\TRule{\True}{\Box}[2]}
                    ]
                  ]
                ]
              ]
            ]
          ]
        ]
      ]
    ]
  \end{oltableau}
  இப்போது வரிகள் 2, 8, 9 ஆகியவற்றில் சரிபார்ப்புக்குறிகள் இல்லை.
  ஆனால் புதிய முன்னொட்டு எதுவும் சேர்க்கப்படவில்லை; ஆகவே விளைவாக வரும் எல்லாக்
  கிளைகளிலும் (அவை மூடாதவரை) வரிகள்~8,~9 க்கு
  $\TRule{\True}{\lor}$ ஐப் பயன்படுத்துகிறோம்:
  \begin{oltableau}
    [\pFmla{\False}{\Box(p \lor q) \lif (\Box p \lor \Box q)}{1},
      just = \TAss, checked
      [\pFmla{\True}{\Box(p \lor q)}{1},
        just = {\TRule{\False}{\lif}[1]}, checked
        [\pFmla{\False}{\Box p \lor \Box q}{1},
          just = {\TRule{\False}{\lif}[1]}, checked
          [\pFmla{\False}{\Box p}{1},
            just = {\TRule{\False}{\lor}[3]}, checked
            [\pFmla{\False}{\Box q}{1},
              just = {\TRule{\False}{\lor}[3]}, checked
              [\pFmla{\False}{p}{1.1},
                just = {\TRule{\False}{\Box}[4]}, checked
                [\pFmla{\False}{q}{1.2},
                  just = {\TRule{\False}{\Box}[5]}, checked
                  [\pFmla{\True}{p \lor q}{1.1},
                    just = {\TRule{\True}{\Box}[2]}, checked
                    [\pFmla{\True}{p \lor q}{1.2},
                      just = {\TRule{\True}{\Box}[2]}, checked
                      [\pFmla{\True}{p}{1.1},
                        just = {\TRule{\True}{\lor}[8]}, checked, close
                      ]
                      [\pFmla{\True}{q}{1.1},
                        just = {\TRule{\True}{\lor}[8]}, checked
                        [\pFmla{\True}{p}{1.2},
                          just = {\TRule{\True}{\lor}[9]}, checked]
                        [\pFmla{\True}{q}{1.2},
                          just = {\TRule{\True}{\lor}[9]}, checked, close]
                      ]
                    ]
                  ]
                ]
              ]
            ]
          ]
        ]
      ]
    ]
  \end{oltableau}
  திறந்த கிளை ஒன்று மீதமுள்ளது; அது முழுமையானது. அதிலிருந்து
  $W = \{1, 1.1, 1.2\}$ என்ற உலகங்களையும் (திறந்த கிளையில் இடம்பெறும்
  முன்னொட்டுகள் இவை மட்டுமே), $R = \{\tuple{1, 1.1}, \tuple{1, 1.2}\}$
  என்ற அணுகுறவையும், $V(p) = \{1.2\}$ (ஏனெனில் வரி~11 ஆனது
  $\sFmla{\True}{p}[1.2]$ ஐக் கொண்டுள்ளது), $V(q) = \{1.1\}$ (ஏனெனில்
  வரி~10 ஆனது $\sFmla{\True}{q}[1.1]$ ஐக் கொண்டுள்ளது) என்ற
  ஒதுக்கீட்டையும் கொண்ட மாதிரியை வரையறுக்கிறோம். மாதிரி
  \olref{fig:counter-Box} இல் படமாக்கப்பட்டுள்ளது; அது $\Box(p \lor q)
  \lif (\Box p \lor \Box q)$ க்கான எதிர்மாதிரி என்பதை நீங்கள்
  சரிபார்க்கலாம்.
  \begin{figure}
  \begin{center}
    \begin{tikzpicture}[modal]
      \node[world] (w1) [label={[align=right]right:\mFalse{p}\\ \mFalse{q}}]{$1$};
      \node[world] (w2) [label={[align=right]right:\mFalse{p}\\ \mTrue{q}},
        above left=of w1]{$1.1$};
      \node[world] (w3) [label={[align=right]right:\mTrue{p}\\ \mFalse{q}},
        above right=of w1] {$1.2$};
      \draw[->] (w1) to (w2);
      \draw[->] (w1) to (w3);
    \end{tikzpicture}
  \end{center}
  \caption{$\Box(p \lor q) \lif (\Box p \lor \Box q)$ க்கான
    ஓர் எதிர்மாதிரி.}
  \ollabel{fig:counter-Box}
\end{figure}
\end{ex}
}
{
\begin{ex}
  $\Proves/ (\Diamond p \land \Diamond q) \lif \Diamond(p \land q)$
  என்று அறிவோம். ஓர் அட்டவணை மரத்தின் கட்டமைப்பு பின்வருமாறு தொடங்குகிறது:
  \begin{oltableau}
    [\pFmla{\False}{(\Diamond p \land \Diamond q) \lif \Diamond(p \land q)}{1},
      just = \TAss, checked
      [\pFmla{\True}{\Diamond p \land \Diamond q}{1},
        just = {\TRule{\False}{\lif}[1]}, checked
        [\pFmla{\False}{\Diamond(p \land q)}{1},
          just = {\TRule{\False}{\lif}[1]}
          [\pFmla{\True}{\Diamond p}{1},
            just = {\TRule{\True}{\land}[2]}, checked
            [\pFmla{\True}{\Diamond q}{1},
              just = {\TRule{\True}{\land}[2]}, checked
              [\pFmla{\True}{p}{1.1},
                just = {\TRule{\True}{\Diamond}[4]}, checked
                [\pFmla{\True}{q}{1.2},
                  just = {\TRule{\True}{\Diamond}[5]}, checked
                ]
              ]
            ]
          ]
        ]
      ]
    ]
  \end{oltableau}
  !!{tableau} இன்னும் முடிவடையவில்லை என்பது தெளிவு. அடுத்த படியில்,
  சரிபார்ப்புக்குறி இல்லாத ஒரே வரியை, அதாவது வரி~$3$ இலுள்ள முன்னொட்டுள்ள
  % SOURCE CORRECTION TA-NML-055: line 3 is false Diamond, expanded by the false-Diamond rule.
  !!{formula}~$\sFmla{\False}{\Diamond(p \land q)}[1]$ ஐக் கருதுகிறோம்.
  % SOURCE CORRECTION TA-NML-058: the search constructs a complete tableau, not a closed one.
  முழுமையான அட்டவணை மரத்தின் கட்டமைப்பின்படி, கிளையில் பயன்படுத்தப்பட்ட
  ஒவ்வொரு முன்னொட்டிற்கும், அதாவது $1.1$, $1.2$ ஆகிய இரண்டிற்கும்,
  $\TRule{\False}{\Diamond}$ விதியைப் பயன்படுத்த வேண்டும்:
  \begin{oltableau}
    % SOURCE CORRECTION TA-NML-056: retain the original implication direction at the root.
    [\pFmla{\False}{(\Diamond p \land \Diamond q) \lif \Diamond(p \land q)}{1},
      just = \TAss, checked
      [\pFmla{\True}{\Diamond p \land \Diamond q}{1},
        just = {\TRule{\False}{\lif}[1]}, checked
        [\pFmla{\False}{\Diamond (p \land q)}{1},
          just = {\TRule{\False}{\lif}[1]}
          [\pFmla{\True}{\Diamond p}{1},
            just = {\TRule{\True}{\land}[2]}, checked
            [\pFmla{\True}{\Diamond q}{1},
              just = {\TRule{\True}{\land}[2]}, checked
              [\pFmla{\True}{p}{1.1},
                just = {\TRule{\True}{\Diamond}[4]}, checked
                [\pFmla{\True}{q}{1.2},
                  just = {\TRule{\True}{\Diamond}[5]}, checked
                  [\pFmla{\False}{p \land q}{1.1},
                    just = {\TRule{\False}{\Diamond}[3]}
                    [\pFmla{\False}{p \land q}{1.2},
                      just = {\TRule{\False}{\Diamond}[3]}
                    ]
                  ]
                ]
              ]
            ]
          ]
        ]
      ]
    ]
  \end{oltableau}
  இப்போது வரிகள் 3, 8, 9 ஆகியவற்றில் சரிபார்ப்புக்குறிகள் இல்லை.
  ஆனால் புதிய முன்னொட்டு எதுவும் சேர்க்கப்படவில்லை; ஆகவே விளைவாக வரும் எல்லாக்
  கிளைகளிலும் (அவை மூடாதவரை) வரிகள்~8,~9 க்கு
  $\TRule{\False}{\land}$ ஐப் பயன்படுத்துகிறோம்:
  \begin{oltableau}
    [\pFmla{\False}{(\Diamond p \land \Diamond q) \lif \Diamond(p \land q)}{1},
      just = \TAss, checked
      [\pFmla{\True}{\Diamond p \land \Diamond q}{1},
        just = {\TRule{\False}{\lif}[1]}, checked
        [\pFmla{\False}{\Diamond(p \land q)}{1},
          just = {\TRule{\False}{\lif}[1]}
          [\pFmla{\True}{\Diamond p}{1},
            just = {\TRule{\True}{\land}[2]}, checked
            [\pFmla{\True}{\Diamond q}{1},
              just = {\TRule{\True}{\land}[2]}, checked
              [\pFmla{\True}{p}{1.1},
                just = {\TRule{\True}{\Diamond}[4]}, checked
                [\pFmla{\True}{q}{1.2},
                  just = {\TRule{\True}{\Diamond}[5]}, checked
                  [\pFmla{\False}{p \land q}{1.1},
                    just = {\TRule{\False}{\Diamond}[3]}, checked
                    [\pFmla{\False}{p \land q}{1.2},
                      just = {\TRule{\False}{\Diamond}[3]}, checked
                      [\pFmla{\False}{p}{1.1},
                        just = {\TRule{\False}{\land}[8]}, checked, close
                      ]
                      [\pFmla{\False}{q}{1.1},
                        just = {\TRule{\False}{\land}[8]}, checked
                        [\pFmla{\False}{p}{1.2},
                          just = {\TRule{\False}{\land}[9]}, checked]
                        [\pFmla{\False}{q}{1.2},
                          just = {\TRule{\False}{\land}[9]}, checked, close]
                      ]
                    ]
                  ]
                ]
              ]
            ]
          ]
        ]
      ]
    ]
  \end{oltableau}
  திறந்த கிளை ஒன்று மீதமுள்ளது; அது முழுமையானது. அதிலிருந்து
  $W = \{1, 1.1, 1.2\}$ என்ற உலகங்களையும் (திறந்த கிளையில் இடம்பெறும்
  முன்னொட்டுகள் இவை மட்டுமே), $R = \{\tuple{1, 1.1}, \tuple{1, 1.2}\}$
  என்ற அணுகுறவையும், $V(p) = \{1.1\}$ (ஏனெனில் வரி~6 ஆனது
  $\sFmla{\True}{p}[1.1]$ ஐக் கொண்டுள்ளது), $V(q) = \{1.2\}$ (ஏனெனில்
  வரி~7 ஆனது
  % SOURCE CORRECTION TA-NML-057: line 7 contains q at prefix 1.2.
  $\sFmla{\True}{q}[1.2]$ ஐக் கொண்டுள்ளது) என்ற ஒதுக்கீட்டையும் கொண்ட
  மாதிரியை வரையறுக்கிறோம். மாதிரி \olref{fig:counter-Diamond} இல்
  படமாக்கப்பட்டுள்ளது; அது $(\Diamond p \land \Diamond q) \lif
  \Diamond (p \land q)$ க்கான எதிர்மாதிரி என்பதை நீங்கள் சரிபார்க்கலாம்.
  \begin{figure}
  \begin{center}
    \begin{tikzpicture}[modal]
      \node[world] (w1) [label={[align=right]right:\mFalse{p}\\ \mFalse{q}}]{$1$};
      \node[world] (w2) [label={[align=right]right:\mTrue{p}\\ \mFalse{q}},
        above left=of w1]{$1.1$};
      \node[world] (w3) [label={[align=right]right:\mFalse{p}\\ \mTrue{q}},
        above right=of w1] {$1.2$};
      \draw[->] (w1) to (w2);
      \draw[->] (w1) to (w3);
    \end{tikzpicture}
  \end{center}
  \caption{$(\Diamond p \land \Diamond q) \lif
    \Diamond (p \land q)$ க்கான ஓர் எதிர்மாதிரி.}
  \ollabel{fig:counter-Diamond}
\end{figure}
\end{ex}
}

\end{document}
